\section{Introduction}
Reusability is one of the main principles in the Knowledge Graph (KG) development process defined by iTelos. The KG project documentation plays an important role in enhancing the reusability of the resources handled and produced during the process. A clear description of the resources, the process (and sub-processes) developed and the evaluation at each step of the process provides a clear understanding of the project, thus serving such information to external readers for the future exploitation of the project's outcomes.\\

\noindent The current document provides a detailed report of the \textit{Ski Resorts in Trentino} project, developed following the iTelos methodology as part of the Knowledge Graphs course (KG 2025/2026) at the University of Trento, under the supervision of Prof. Fausto Giunchiglia and the KnowDive Research Group. The report is structured as follows:
\begin{itemize}
    \item Section 2: Definition of the project's purpose (Domain of Interest, informal purpose) and related information gathering.

    \item Sections 3, 4, 5, 6: The description of the iTelos process phases (Purpose Formalization, Information Gathering, Language Definition, Knowledge Definition, Entity Definition) and their activities, divided by knowledge and data layer activities, as well as the evaluation of the resources produced in terms of fitness for the chosen purpose.

    \item Section 7: The description of the metadata produced for all (and all kinds of) resources handled and generated by the iTelos process, while executing the project.

    \item Section 8: Conclusion and open issues summary.
\end{itemize}
